% ══════════════════════════════════════════════════════════════════
% 第一章 绪论
% ══════════════════════════════════════════════════════════════════
\chapter{绪\hspace{6pt}论}

\section{研究背景与意义}

大型工业基础设施——桥梁、压力容器、风力发电机组和石油天然气管道——在其服役周期内持续承受疲劳载荷、腐蚀环境和意外冲击，表面及近表面裂纹、壁厚减薄和材料退化是其最常见的损伤模式。对这类结构实施定期无损检测并非可选维护策略，而是保证公共安全和避免灾难性失效的强制性工程需求\citing{qu2020development}。然而，受检结构动辄数十米高的作业面、受限空间内不可达的检测点、以及高空作业固有的安全风险，使传统人工搭架检测在效率、成本和人员安全三个维度上均面临难以调和的矛盾\citing{ball1998detection}。

多旋翼无人机为上述困境提供了独特的解决方案。与固定翼或大型直升机平台不同，多旋翼的悬停能力、低空机动性和对复杂几何表面的适应性使其成为近距离结构检测的理想空中载体。过去十年间，无人机平台从简单的航拍摄像发展为搭载高清可见光相机、红外热像仪和激光雷达的多传感器系统，实现了从远距离目视筛查向近场定量测量的功能跃迁。然而，视觉与热成像等非接触传感模态共享一个根本局限：它们只能探测表面或近表面（对热波而言）异常，对隐藏在材料内部的亚表面分层、壁厚减薄和体缺陷缺乏穿透能力。

电磁超声换能器（EMAT）的引入打破了上述深度限制。EMAT利用交变线圈在导电材料集肤深度内感生涡流，涡流与永磁体提供的静偏置磁场耦合产生洛伦兹体力，后者直接作为体载荷源在材料内部激发超声波\citing{sun2025emat}。这一换能机理带来了两个关键优势：第一，超声波在工件内部而非探头中产生，消除了传统压电超声对液体耦合剂的依赖，使检测可在干燥、高温和粗糙表面上进行；第二，回波信号携带了完整的传播路径信息——从激励线圈的涡流分布、到材料内部的弹性波传播、再到底部界面的反射——由此可同时提取壁厚、缺陷位置和材料声学属性。在静态实验室条件下，基于EMAT的测厚精度已达到亚毫米级，测量误差低于1\%\citing{kubrusly2023measurement}，且在500°C高温和5 mm提离距离下仍能保持优于3.5\%的精度\citing{chen2024novel}。

然而，EMAT的上述精度承诺建立在探头与被测表面之间保持稳定相对位置的前提之上——这一前提在无人机搭载条件下恰恰被根本性地打破了。问题的核心在于：EMAT的换能效率对探头-表面间距（升距）极其敏感。升距的微小变化会通过两种机制同时恶化信号质量：磁通密度随距离按负幂律衰减，导致涡流密度和洛伦兹力幅值下降；涡流的空间分布也因磁力线扩散而改变，使激励声场的指向性和模态纯度退化。当EMAT探头被刚性安装在多旋翼机体上时，平台自身的姿态微振（±2°量级）、螺旋桨尾流引起的不稳定地面效应、以及接触瞬间的动力学反冲，都会在亚毫米尺度上持续调制升距——而这一调制尺度恰好与EMAT的敏感区间重叠。

更根本的困难在于，上述升距变化并非以孤立方式发生。无人机的法向速度收敛过程、末端执行器的力学刚度、被测表面的曲率和材料属性，以及环境风场的随机扰动，共同构成了一个强耦合的多因素系统。在该系统中，接触状态的改变不是简单的二值事件，而是升距逐渐减小、法向速度渐近收敛、回波能量非线性增长和频谱特征持续演化的连续动态过程。单一传感器——无论是超声回波幅值、飞控位姿残差还是视觉深度估计——都无法独立可靠地表征这一复杂过程。当升距因姿态扰动而瞬时增大时，回波幅值的下降既可能源于真实的接触脱离，也可能来自暂时性的磁耦合减弱，而后者在姿态恢复后会自动消失。依赖单帧幅值阈值的传统判别方法在物理上无法区分这两种场景，其误判不可避免地导致控制闭环中的频繁状态切换——即接触-脱离-再接触的循环振荡。

上述分析指向一个清晰的研究空白：无人机EMAT检测需要一个能够联合解释电磁-弹性耦合、飞行器刚体动力学和视觉几何约束的多模态接触感知框架，而目前的方法论——无论是纯数据驱动的黑箱分类器，还是基于单一物理量阈值的工程判据——均不具备在动态扰动下维护物理一致性的能力。纯数据驱动方法缺乏对EMAT换能链中物理因果关系的显式建模，在传感器噪声增大或环境条件改变时泛化能力急剧下降；物理约束方法（如基于PINN的动力学建模）虽然引入了方程正则化，但大多局限于单一连续系统的状态估计，未涉及跨模态、跨时间尺度的异步观测融合问题。

本文正是在这一研究空白处开展工作。我们提出一种基于多模态逐级物理约束的框架，将EMAT超声回波序列、无人机六自由度位姿残差和RGBD局部视觉特征统一建模为一个时序感知系统。在该框架中，物理先验——包括时间邻近性、空间邻近性和法向一致性——被显式编码为Transformer注意力机制的对数偏置项，而非被隐含在数据分布中。这一设计使模型能够在飞行扰动下维持可解释的证据链：视觉确认几何可达性，位姿验证动力学收敛性，超声评估电磁耦合质量，三者同时满足时接触判定才被确立。该框架的理论价值在于将接触状态从工程经验判据提升为具有物理可证伪性的概率估计问题；其工程价值在于为无人机自主无损检测提供了一条从传感器原始输出到闭环控制决策的完整推理路径，有望推动多旋翼平台从视觉巡检工具向真正的空中物理交互系统的范式转变。

\section{国内外研究现状}

\subsection{非侵入接触式检测无人机研究现状}

结构健康监测领域对空中检测平台的需求源于一个无法调和的工程矛盾：大型工业设施——化学反应器、储罐壁板和蒸汽锅炉——的检测覆盖率要求与人工可达性之间存在根本性鸿沟。搭设脚手架或动用高空作业车不仅成本高昂（单次检测动辄数十万人民币），而且在密闭空间、高耸结构和离岸设施等场景中存在致命安全风险\citing{qu2020development}。

这一矛盾驱动了将多旋翼无人机用作移动传感平台的研究方向，其发展轨迹呈现出一条清晰的技术路线演化：从被动的远距离目视记录，到主动的物理激励与信号采集，再到自主决策与感控闭环。图\ref{fig:industrial_structures}展示了无人机无损检测所面向的三类典型复杂工业结构及其代表性缺陷模式。

\begin{figure}[htbp]
    \centering
    \begin{subfigure}{0.32\textwidth}
        \centering
        \includegraphics[width=\textwidth]{c1/reactor.jpg}
        \caption{}
    \end{subfigure}
    \hfill
    \begin{subfigure}{0.32\textwidth}
        \centering
        \includegraphics[width=\textwidth]{c1/storage_tank.jpg}
        \caption{}
    \end{subfigure}
    \hfill
    \begin{subfigure}{0.32\textwidth}
        \centering
        \includegraphics[width=\textwidth]{c1/steam_boiler.png}
        \caption{}
    \end{subfigure}
    \caption{无人机无损检测面向的典型复杂工业结构。(a) 化学反应器——长期暴露于辐射、高温和腐蚀介质，表面裂纹、分层和前沿侵蚀是常见损伤模式；(b) 储罐壁板——受内压与腐蚀介质作用，壁厚减薄和点蚀是检测重点；(c) 蒸汽锅炉——焊缝退化和应力腐蚀裂纹是主要安全隐患。三类结构均位于高空或受限空间，构成无人机接触式无损检测的核心应用场景。}
    \label{fig:industrial_structures}
\end{figure}

无人机用于结构检测的早期实践——约2010至2015年——本质上是对人工目视检查的替代和延伸。搭载高清可见光相机的四旋翼平台使检测人员能够在地面安全位置获取结构表面的高分辨率图像，而红外热像仪则拓展了对近表面脱粘、水分侵入和隔热层缺陷的探测能力\citing{omar2017uavir}。钟嶒楒等人\citing{zhong2026sam}面向锅炉水冷壁缺陷提出融合SAM分割与YOLO的三阶段检测方法，以实例分割处理常规与极小缺陷、目标检测器识别极大目标、SAM大模型提供掩膜支持，在爬壁机器人采集的水冷壁图像上实现了97.3\%的平均检测准确率（漏检率0.75\%），并可耦合无人机平台推动锅炉检测的自动化与智能化。郑桂红等人\citing{zheng2021boiler}基于电站锅炉与垃圾焚烧锅炉的内部检验案例，总结了面向炉内结构的无人机检验方案制定方法与改进方向，展示了无人机在锅炉内部受限空间检验中的工程可行性。针对锅炉炉膛等受限空间，张冀等人\citing{zhang2015waterwall}提出遥控无人机进入炉膛录像并与数据库图像比对分析的锅炉水冷壁故障检测方法，免除了人工进炉检测的高风险。图\ref{fig:visual_inspection_evolution}展示了从可见光巡检到红外热成像等多模态无损检测的技术演进路径。

\begin{figure}[htbp]
    \centering
    \begin{subfigure}{0.32\textwidth}
        \centering
        \includegraphics[width=\textwidth]{c1/feroz2021uav.png}
        \caption{}
    \end{subfigure}
    \hfill
    \begin{subfigure}{0.32\textwidth}
        \centering
        \includegraphics[width=\textwidth]{c1/memari2024windturbine.png}
        \caption{}
    \end{subfigure}
    \hfill
    \begin{subfigure}{0.32\textwidth}
        \centering
        \includegraphics[width=\textwidth]{c1/omar2017uavir.png}
        \caption{}
    \end{subfigure}
    \caption{无人机多模态视觉巡检。(a) 利用无人机搭载高清可见光相机对桥梁进行表面裂纹检测；(b) 利用无人机搭载红外热像仪对混凝土桥面板进行非接触脱粘检测；(c) 利用无人机搭载可见光相机与激光雷达对风电叶片进行裂纹、侵蚀和涂层异常等缺陷检测。}
    \label{fig:visual_inspection_evolution}
\end{figure}

Nooralishahi等人\citing{nooralishahi2021review}对工业场地无人机无损检测的综述进一步确认，多传感载荷组合已成为该领域的标准配置，同时指出检测可靠性与自主性仍是尚未解决的问题。这一阶段的系统架构相对简单：无人机遥控、数据采集与缺陷判别完全由地面操作人员完成。尽管基于可见光成像的无人机无损检测在表面筛查方面取得了显著进展，其功能边界也十分明确——成像传感器仅能捕获表面或热扩散可达的浅表层异常，对材料内部的分层扩展、壁厚减薄和体缺陷缺乏穿透性检测能力\citing{ball1998detection}。

综上所述，无人机视觉检测凭借部署灵活、覆盖效率高等优势，已在反应器、储罐及锅炉等场景中得到广泛研究。然而，此类方法的检测主要适用于表面及浅表层异常，难以对材料内部进行定量评估。因此，实现从表面筛查到内部检测的拓展，仍需引入超声、电磁涡流等主动检测手段。沿这一方向，研究者先后发展了壁面吸附式超声检测无人机\citing{mattar2018wallsticking}、自主空中超声测厚系统\citing{zhang2021autonomous}与无人机冲击回波混凝土检测系统\citing{thurnherr2025impactecho}；近期，Veenstra等人\citing{veenstra2026autonomous}更在商用四旋翼平台上实现了非结构化环境中的全自主接触式超声检测，标志着无人机检测正从目视筛查向主动声学检测稳步过渡。

金属结构壁厚的定量测量——无论是管道腐蚀评估、储罐剩余寿命预测还是压力容器安全评级——是工业无损检测中最基础也最频繁执行的任务之一。厚度测量的准确性直接决定了结构完整性判定的可信度：对于壁厚减薄型损伤，亚毫米级的测量偏差可能意味着服役寿命预测中数年量级的误差。这一精度要求驱动了超声测厚技术从压电耦合到电磁非接触的持续演进。EMAT测厚的现有研究几乎全部在静态或准静态条件下完成：探头固定在精密位移台上，升距由微米级螺旋千分尺设定，被测表面为平面且表面状态已知，信号采集具有充足的时间进行多次平均以提高信噪比\citing{kubrusly2023measurement,chen2024novel}。上述受控条件排除了无人机搭载场景中最核心的困难来源——平台运动引起的时变边界条件。在无人机飞行检测中，即使是最稳定的悬停状态也伴随±2°的姿态波动和毫米级的低频位置漂移，这些扰动通过前文分析的升距非线性直接转化为接收信号的幅值调制和相位抖动。当回波到达时刻因信噪比不足而无法可靠提取时——EMAT的典型信噪比仅为15-30 dB，远低于压电超声的40-60 dB——TOF估计的方差将急剧增大，直接导致厚度反演精度的崩溃\citing{sun2025emat}。

更隐蔽的困难在于波传播路径的几何畸变。在曲面试件上，底部反射面的曲率使回波能量并非集中返回探头位置而是向空间弥散，降低了有效接收能量；当探头法向与表面局部法向存在偏差角时，激励声束的主瓣偏离底部反射面的法线方向，使往返路径不再满足$d = c\Delta t / 2$的简单几何关系——精确反演需要已知表面曲率和入射角，而这两个信息在无人机检测中恰恰无法以足够精度获取。曲面、运动、低信噪比和非理想耦合四个因素的叠加，使得动态工况下的EMAT测厚成为一个远超静态实验复杂度的多物理场逆问题。国内已出现面向无人机平台的电磁超声测厚装置与系统设计\citing{wu2021ematdrone,lin2026furnace,zheng2026anticollision}，但这些工作集中于装置集成与检测流程层面，尚未触及动态工况下的信号退化机制与反演精度问题。

当前研究在处理上述问题时存在两个根本性不足。其一，信号处理策略几乎全部建立在对单帧数据的独立分析之上，缺乏利用多帧时序信息进行联合估计的能力——而无人机的运动学状态（可以通过IMU和里程计以100 Hz以上的频率获取）恰为相邻帧之间的厚度变化提供了强先验约束：无人机在接触维持期间的法向位置变化通常不超过每秒数毫米，因此相邻帧的厚度估计不应出现跳跃式变化。其二，现有方法缺乏将EMAT换能的物理模型——包括升距-涡流-信号幅值的映射关系和波速-温度-应力的耦合关系——作为显式约束嵌入厚度反演算法的机制。纯数据驱动的反演方法在训练数据覆盖的工况范围内可表现良好，但一旦升距范围、表面曲率或材料类型超出训练分布，其精度和可靠性便无法保证。这两个不足共同指向一个清晰的研究方向：将无人机运动学状态的时序约束和EMAT换能机理的物理约束联合嵌入厚度反演框架，使反演过程不仅受到数据的驱动，更受到物理可解释性的规范——这正是本文逐级物理约束方法论在电磁超声测厚维度的核心切入点。

上述讨论从EMAT信号本身的物理特性出发，分析了动态工况下测厚精度退化的机制根源。然而，EMAT检测还有一个更为前端的维度尚未被充分审视：与远距离视觉成像不同，超声、电磁涡流等主动检测手段要求传感器与被测表面保持稳定的近场距离或物理接触，使无人机面临气动稳定、接触力控制和末端定位等方面的挑战。在EMAT探头接触被测表面之前，无人机平台自身如何完成精确且稳定的贴靠？这一问题将分析视角从“检测信号的质量”拉回到“产生信号的前提条件”——即空中物理交互的稳定性，也构成了下文精确接触感控研究的重要背景。

\subsection{非侵入接触式检测无人机精确接触感控研究现状}

无人机无损检测系统的演化遵循一条由传感器模态驱动的技术路径。当采集数据的传感器从被动接收型（相机、热像仪）转变为主动激励-接收型（超声探头、涡流线圈、EMAT）时，无人机平台的物理角色也随之发生了根本性转变——从在安全距离外静默观察的飞行相机平台，转变为必须与被测表面建立并维持近场物理耦合的空中交互系统。这一转变引入了一类全新的工程问题，其本质可归结为空中物理交互（Aerial Physical Interaction, APhI）的核心矛盾：如何在气动载荷持续扰动、机体动力学强非线性、且执行器数目受限的条件下，实现末端效应器与外部结构之间精确、稳定且可重复的物理接触。Ollero等人\citing{ollero2021past}对空中机械手（aerial manipulator）过去、现状与未来的系统综述确认了这一领域已从概念验证走向工程应用，其中工业检测与维护是主要落地场景之一；从AEROARMS项目的多自由度空中机械臂\citing{ollero2018aeroarms}到面向油气工业的接触检测空中机械手\citing{trujillo2019novel}，接触式无损检测始终是推动空中机械手发展的重要需求。

标准四旋翼系统属于欠驱动系统（underactuated system）：四个旋翼推力控制了六个空间自由度，水平面内的平动必须通过机体倾斜产生的推力水平分量来实现。这一约束在接触式检测中制造了一个根本性的运动学冲突。当探头需要向竖直壁面施加法向接触力时，机体必须向壁面方向倾斜以获得水平推力分量，而这一倾斜角恰使探头轴线偏离壁面法线方向。换言之，推力方向的建立以牺牲姿态对准为代价——在同一控制周期内，飞行器无法同时满足法向力施加（需要倾斜）和探头法向对准（需要垂直）两个互斥的目标。对于需精确维持升距和法向角的EMAT检测，这一冲突的影响远比压电超声严重：压电探头在干耦合弹性体的柔顺范围内可容忍数度的姿态偏差，而EMAT的探头-表面角度偏差超过1-2°即开始显著改变涡流空间分布和激励声场指向性。

应对这一冲突的第一条技术路线是被动柔顺机构设计。接触式超声检测的空中实施面临的首要困难来自接触瞬间的动力学冲击：多旋翼无人机在自由飞行阶段依靠推力矢量维持姿态和位置，当末端探头突然与被测表面发生物理接触时，接触反力会瞬间打破力-力矩平衡——这一瞬态扰动的幅值可达数牛顿，持续时间数十毫秒\citing{gonzalez2019payload}。若不加专门处理，飞控系统会将接触力视为外部位置扰动并试图通过增大推力来补偿，反而导致探头与表面之间的接触力失控增长，造成检测失败甚至平台失稳。Gonz\'{a}lez-deSantos等人\citing{gonzalez2019payload}通过在探头与机体之间嵌入弹簧-阻尼串联元件，将接触冲击动能转化为弹性势能和热耗散，有效降低了传递至飞控回路的扰动幅值。同一思路的变体包括Hamaza等人\citing{hamaza2018adaptive}提出的自适应柔顺空中机械臂，通过在线调节控制器增益定制末端施力特性；Gonz\'{a}lez-deSantos等人\citing{gonzalez2020collision}进一步设计了带碰撞缓解机构的检测载荷，使无人机在无GNSS辅助的受限环境中仍能维持与垂直结构的稳定接触。此类方案结构简单、响应速度快（机械元件的力学响应滞后在毫秒量级），无需传感器反馈和计算资源，适用于平面或小曲率表面上的初始接触建立。然而，被动柔顺的局限性同样突出。弹簧的刚度-行程乘积决定了可吸收能量的上限，而这一乘积受探头质量和可用安装空间的双重约束；弹簧恢复力还会在接触维持阶段持续施加一个与形变成正比的推力，使得真正的法向接触力无法被独立设定——力的大小由弹簧预压缩量和当前相对位移决定，而相对位移又受飞控定位精度的制约。在亚厘米级定位误差的典型条件下，高刚度弹簧对应的接触力变化可达数牛顿，远超过EMAT稳定工作区间（<1 N的力波动即可引起回波幅值显著变化）。当检测要求恒定的法向接触力以维持稳定的声耦合条件时——例如压电探头需要的10-20 N等效压紧力——被动柔顺便力不从心。

第二条技术路线——主动力控——通过将接触力纳入无人机控制回路的反馈输入，实现了对法向力的闭环调节。Kocer等人\citing{kocer2019inspection}利用非线性模型预测控制（NMPC）将未知接触扰动作为外部力进行在线估计和补偿，成功在5.5 N恒定接触力下获取了天花板结构的超声A扫描数据。主动力控方案的关键优势在于其对不同表面刚度和几何形状的自适应能力：控制器可以根据力传感器的反馈动态调整推力输出，无需在检测前获知准确的表面机械阻抗。然而，NMPC的在线优化计算在高频飞控循环（100 Hz以上）中的实时性保证依赖于简化的接触动力学模型，而简化的代价是对未建模模态的敏感性增加。当无人机在真实环境中遭遇阵风扰动、螺旋桨尾流的地面效应或被测结构的力学阻抗突变时，预测模型与物理系统之间的偏差可能导致力跟踪性能的瞬时退化——而在EMAT检测中，即使数百毫秒的升距失控也足以导致有效回波数据的完全丢失。图\ref{fig:contact_compliance_vs_force}对比了被动柔顺机构与主动力控NMPC框架的典型实现。

\begin{figure}[htbp]
    \centering
    \begin{subfigure}{0.48\textwidth}
        \centering
        \includegraphics[width=\textwidth]{c1/gonzalez2019payload.png}
        \caption{}
    \end{subfigure}
    \hfill
    \begin{subfigure}{0.48\textwidth}
        \centering
        \includegraphics[width=\textwidth]{c1/kocer2019inspection.png}
        \caption{}
    \end{subfigure}
    \caption{接触式超声检测的两条典型技术路线。(a) Gonz\'{a}lez-deSantos等人\cite{gonzalez2019payload}提出的被动柔顺末端执行器——弹簧-阻尼串联机构吸收接触冲击动能，但恢复力$F_n = k\cdot\Delta x$导致法向力与形变不可控耦合；(b) Kocer等人\cite{kocer2019inspection}提出的基于NMPC/NMHE的主动力控框架——将未知接触扰动作为外部力在线估计并补偿，实现法向力的闭环独立调节。}
    \label{fig:contact_compliance_vs_force}
\end{figure}

第三条技术路线通过增加独立推力执行器的数量——可倾转旋翼、矢量推力单元或额外的侧向推进器——将标准欠驱动四旋翼升级为全驱动（fully-actuated）和过驱动（over-actuated）平台，使推力矢量的可控维数扩展至6，实现位置控制与姿态控制的动力学解耦。在此类平台上，法向力的施加不再需要机体倾斜，探头可与壁面始终保持法向对准。Watson等人\citing{watson2022dry}利用Voliro全矢量推力六旋翼完成了倾斜表面和悬垂面上的干耦合超声B扫描，法向力维持在约20 N；Marcellini等人\citing{marcellini2024development}在倾转旋翼平台上集成了双重阻抗控制，实现机械臂端精细力控与机体端姿态稳定的解耦运行；Rashad等人\citing{rashad2020towards}进一步将视觉信息引入阻抗控制框架，在线估计未知曲面的局部法向与几何形状，将末端执行器期望轨迹投影至被测表面，在无需先验几何模型的条件下完成贴靠接触。图\ref{fig:overactuated_platforms}展示了全驱动平台在接触式无损检测中的两种典型实现。

\begin{figure}[htbp]
    \centering
    \begin{subfigure}{0.48\textwidth}
        \centering
        \includegraphics[width=\textwidth]{c1/watson2022dry.png}
        \caption{}
    \end{subfigure}
    \hfill
    \begin{subfigure}{0.48\textwidth}
        \centering
        \includegraphics[width=\textwidth]{c1/marcellini2024development.png}
        \caption{}
    \end{subfigure}
    \caption{全驱动平台在接触式无损检测中的两种典型实现。(a) 利用Voliro全矢量推力六旋翼——通过可倾转旋翼实现位置-姿态动力学解耦，在倾斜和悬垂表面上完成干耦合超声B扫描，法向力维持在~20 N；(b) 基于倾转旋翼平台的半自主无损检测框架——双重阻抗控制（机械臂端精细力控 + 机体端姿态稳定）解耦运行，动态吸收接触冲击扰动。}
    \label{fig:overactuated_platforms}
\end{figure}

全驱动方案的根本代价在于机械复杂度和能效的恶化。可倾转机构引入了旋转关节——额外的潜在故障点，在工业现场的粉尘、湿度和腐蚀性环境中可靠性退化速率远高于固定旋翼。倾转执行器的响应带宽也限制了力控环路可达到的动态性能上限：旋翼倾转的角加速度受限于舵机扭矩和旋翼转动惯量，在高频扰动（如阵风引起的机体横滚，频率可达2-5 Hz）下可能无法充分跟踪期望的力输出。此外，执行器数量的增加直接推高了悬停功率需求，缩短了有效检测航时——对于需要大面积扫查的工业检测任务，这是一个不可忽视的操作约束。

上述被动柔顺、主动力控与全驱动三条技术路线均从机械运动学与动力学角度追求物理接触的建立与维持，而几乎未触及一个更深层的问题：物理接触是否等同于有效的超声信号耦合？

当压电超声探头通过干耦合弹性体与金属表面接触时，只要接触正压力超过一定阈值（通常为10-20 N），弹性体即可有效填充微观气隙，实现可重复的声阻抗匹配。然而，电磁超声换能器的耦合机制完全不同：EMAT的换能效率取决于探头线圈与被测表面之间的电磁场相互作用，而这一作用对线圈-表面间距——即升距——具有指数级的敏感性。升距每增加1 mm，涡流密度可下降30-50\%，而涡流密度的下降直接导致洛伦兹体力幅值的等比例衰减。

这一物理差异带来了一项无人机平台特有的工程困境。在压电检测中，接触力是声耦合的直接指标——力越大，耦合越充分，两者的关系单调且在相当宽的力区间内保持近线性。而在EMAT检测中，接触力本身并非耦合质量的可靠指标：即使法向力稳定，微小的姿态角扰动（±1°量级）也会通过改变线圈与表面的相对倾斜角而影响涡流的空间分布和磁力线的穿透深度。换言之，压电检测只需要维持一个标量（接触力），而EMAT检测需要同时维持一个包括升距、法向角和切向滑移在内的三维相对位姿约束。前者可由单输入单输出的力控环处理，后者则需要以多模态感知为基础的联合状态估计——这就将问题从“如何建立并维持接触”推进到“如何感知接触的质量并据此调节平台状态”，即下文反馈微调感控研究要回答的问题。

\subsection{非侵入接触式检测无人机反馈微调感控研究现状}

将上述三条技术路线——欠驱动约束、被动/主动柔顺与全驱动方案——各自的内在局限置于无人机自主检测的完整操作链中，可辨识出一个更深层的系统架构缺陷。当前的系统设计将三个核心功能——感知（传感器数据获取与处理）、控制（飞行动力学调节）和检测（无损检测信号生成与判读）——组织为三个串行独立的子系统：感知为控制提供位姿参考，控制将平台移动至指定位置，检测在位置到达后被动运行。这一串行架构隐含的假设是，各子系统的运行时间尺度可清晰分离——感知和控制快于检测需求，检测不影响感知和控制的运行策略。

在无人机EMAT检测的真实物理系统中，这一假设被彻底打破。EMAT信号质量对升距的亚毫米敏感性意味着，即使在飞控闭环维持的“悬停”状态下，残余的姿态微振（±2°）和位置漂移仍在持续调制超声回波幅值和相位。换言之，控制子系统“认为”已达到稳态的行为，在检测子系统的物理量尺度上仍是持续的瞬态扰动。反过来，检测子系统输出的信号质量信息——回波能量趋势、信噪比变化、频谱重心偏移——本可作为控制子系统的附加反馈，用于在信号质量下降时主动调整贴靠策略，但当前架构缺乏将检测质量信息反馈至控制决策的闭环通道。现有的视觉伺服方案\citing{rashad2020towards}提供了基于图像的位姿反馈，但未融合超声信号质量信息；现有的超声检测系统被动接收飞控输出的位姿带来的信号波动，但未主动引导飞行器调整至最优信号获取姿态。

在已有工作中，信号质量反馈至控制回路的闭环仅以零散形式存在。基于接触力的反馈调节最为成熟——Kocer等人\citing{kocer2019inspection}与Marcellini等人\citing{marcellini2024development}将力传感器读数作为反馈信号调节推力或机械臂力矩，使接触力维持在目标区间；但如上一小节所述，接触力稳定并不保证电磁耦合质量，EMAT场景下的升距与法向角误差无法从力读数中解析。姿态对准则被作为接触建立前的开环几何条件处理——Watson等人\citing{watson2022dry}的干耦合轮式探头要求探头轴线保持在表面法向±10°以内以获得有效声耦合，但该容差依靠平台的全驱动解耦能力在接触前一次性满足，接触维持期间不存在依据回波质量在线修正姿态的机制。回波信号质量通常仅作为检测数据有效性的判据（筛选可用帧、剔除噪声帧），而非控制回路的输入。与本文场景最接近的是Yung等人\citing{yung2022visually}提出的EMAT视觉编码接触检测系统，其通过机器人搭载的立体视觉获取探头三维位姿并融合检测数据生成三维地图，但反馈依据是视觉几何而非超声回波质量，且平台限于地面机器人。将回波特征——能量、信噪比、到达时间稳定性与厚度估计重复性——反馈至飞行控制，驱动横滚/俯仰方向的姿态微调以改善电磁耦合，这一闭环在现有文献中尚属空白。该空白恰与EMAT换能链对升距和法向角的超高灵敏度形成尖锐矛盾：姿态微调所能纠正的正是对EMAT耦合最致命的扰动量，而缺乏回波质量反馈的悬停控制恰恰无法感知这些扰动。

这种感知-控制-检测之间的单向串行关系——而非双向协同闭环——构成了无人机自主EMAT检测中尚未被学界系统论述的关键系统级瓶颈。突破这一瓶颈需要两个能力：一是将多模态感知信息（视觉几何、位姿残差、超声回波特征）在统一时序框架内融合为对接触状态的联合概率估计；二是将该估计以物理可解释的置信度形式输出，使其可直接作为控制状态机的切换依据。这两个能力恰构成了本文逐级物理约束方法论的两个核心支柱：物理约束注意力（实现多模态信息的物理一致融合）和双阈值持续时间滞回判据（实现从感知概率到控制决策的平滑映射）。

\subsection{多模态物理信息Transformer研究现状}

前述分析从平台力学和信号物理两个维度分别揭示了各自的关键瓶颈——升距敏感性导致的力学-电磁二元分离、动态工况下的厚度反演退化、以及回波质量反馈链路的缺失。但这些瓶颈在算法层面共享一个深层原因：多模态信息（视觉、位姿、超声）的融合缺乏物理一致性的保障机制。飞行控制的反馈回路与检测信号的反馈回路本质上是异步、异质且缺乏统一解释框架的：位姿控制器以100 Hz以上的频率对位置和速度误差做出响应，其优化目标是使六自由度状态逼近给定参考值；而EMAT信号质量受升距波动、涡流分布变化和材料电磁属性共同调制，其退化模式无法单纯从位姿误差中推导。当前的主流做法是将控制回路和检测回路视为两个独立子系统——前者追求飞行稳定，后者被动接受信号波动——而缺少一个将两者耦合起来的中间层：一个能够联合解释位姿动力学、电磁-弹性耦合和视觉几何约束的多模态感知模型。以下从多模态融合这一算法视角出发，追溯从统计融合到物理约束嵌入的范式演化，并指出当前方法在物理生成机制约束方面的结构性缺失。

多模态感知信息的融合——将EMAT超声回波、无人机六自由度位姿残差、RGBD视觉深度与纹理特征整合为对接触状态的统一判断——是无人机自主电磁超声检测在算法层面的核心挑战。这一挑战的本质并非简单的“多种数据源的拼接”，而是如何在物理生成机制根本不同的多个观测通道之间建立一致的、可解释的联合表征。围绕这一问题的研究范式经历了从手工统计融合到深度学习表征学习、再到Transformer全局建模\citing{xie2026mpfusionnet}和物理约束嵌入的持续演化，而每一次范式转移都源于对前一阶段方法在特定物理条件下失效模式的识别与回应。

多模态融合的早期方法——包括特征级联（early fusion）、决策级加权（late fusion）和基于卡尔曼滤波的线性高斯融合——均建立在一个共同的统计假设之上：各模态观测噪声相互独立且服从高斯分布，融合权重的确定仅依赖于各通道的估计协方差。当模态数量少（2-3个）、观测模型为线性、且各传感器的物理生成过程保持平稳时，这类方法在理论和工程上都是充分且可证明最优的。然而，无人机EMAT检测恰好系统性地违背了上述三个条件。EMAT的升距-信号幅值关系是高度非线性的（负幂律衰减加磁力线扩散效应），位姿误差的协方差在接触瞬间因冲击扰动而呈现强时变性，视觉深度估计的噪声则不满足高斯假设——在反光、纹理贫乏和遮挡区域表现为具有厚尾分布的离群噪声。这三个条件的共同失效使得基于解析方差的统计融合权重在关键过渡段给出系统性的错误置信度分配。

深度学习表征融合方法——通过CNN、LSTM或MLP将各模态分别编码为隐向量后联合优化——以端到端学习的策略绕过了噪声分布假设的显式建模需求，在静态测试集上取得了远超统计方法的性能。然而，表征融合的根本局限在于其隐空间中缺乏对模态间物理因果关系的编码。当训练数据的分布覆盖了所有工况组合时，表征学习可隐式地捕获模态间的相关性模式；但当测试工况涉及训练分布外的升距范围、表面曲率或环境光照条件时，隐空间中习得的模态关联——其本质只是数据中共现模式的统计拟合——可能产生物理上无法解释甚至矛盾的推断。数据驱动的多模态融合方法（如基于Transformer的序列模型）在大规模语义任务中表现出色，但其注意力权重完全由特征相似度驱动，在面对物理测量噪声和传感器退化时可能产生不可解释甚至物理上矛盾的跨模态关联。

Transformer架构通过多头自注意力和交叉注意力机制实现了对任意位置令牌对之间依赖关系的全局建模，这一能力在自然语言和计算机视觉任务中已被充分验证\citing{vaswani2017attention}。将Transformer引入无人机多模态融合的动机十分明确：EMAT回波的当前帧与数百毫秒前的姿态加速度可能存在长程物理因果关系（例如，一次姿态突然修正引起的升距变化需经过声波往返材料内部后才能在回波中体现），而CNN的感受野或LSTM的门控记忆衰减都可能丢失这一跨时间尺度的依赖。

然而，Transformer在处理物理测量信号时暴露出一种根本性的机制缺陷，其根源在于注意力权重的本质含义：查询向量$\mathbf{Q}$与键向量$\mathbf{K}$的点积衡量的是两个令牌在习得嵌入空间中的统计相似性，而非物理因果性。在自然语言任务中，词嵌入空间中的相似性通常与语义相关性高度一致——两个词在嵌入空间中距离近，意味着它们在大量语料中经常共现，而这种共现往往反映了真实的语义关系。但在物理测量系统中，嵌入空间中的相似性并不具备等价的物理因果含义。EMAT回波能量与视觉深度梯度在嵌入空间中可能因升距变化这一共同隐变量而产生统计相关性——两者都随升距改变而单调变化——但这一相关性的方向、强度和函数形式完全取决于材料属性、磁体几何和工作频率等物理参数，而非任何可推广的语义规律。当这些物理参数偏离训练条件时，习得的嵌入空间相关性可能以虚假的高注意力权重将升距波动引起的视觉噪声错误地解释为超声信号的增强或衰减。

此外，无人机动态环境中有三个特征加剧了上述因果-统计不一致的风险。第一，强烈的非平稳性：升距、姿态角和气流条件在接触建立与脱离过程中持续变化，使各模态特征的边际分布和联合分布不断漂移，而Transformer的注意力矩阵没有内置的分布漂移检测或补偿机制。第二，异步物理延迟：视觉曝光（~15 ms）、EMAT USB传输（~5 ms）和飞控EKF滤波（~10 ms）之间存在不可忽略且非恒定的延迟差异，导致来自同一物理事件的多模态观测在Token序列中被分配至不同时间位置，从而在注意力计算中产生时序错位。第三，非均匀噪声污染：在视觉反光、EMAT电磁干扰或IMU振动饱和等退化工况下，某一模态的信噪比可能在单帧内急剧下降数十dB，而Transformer没有内置机制来识别并降权这种瞬时退化的令牌——相反，高能量的噪声成分可能在注意力计算中反而获得了过高的权重。

为克服纯数据驱动Transformer在物理系统中的上述局限，研究者发展了两类将物理先验嵌入深度学习模型的策略。

第一类——物理信息神经网络（Physics-Informed Neural Networks, PINN）——通过自动微分技术将支配系统行为的偏微分方程（如Navier-Stokes方程、弹性动力学方程或刚体运动学方程）作为残差正则化项附加于损失函数\citing{sanyal2023ramp,gu2024physics}。PINN的物理残差损失在解的可行域内施加了强全局约束，能够有效过滤违背物理定律的输出。然而，PINN的成功高度依赖于对支配方程形式的精确已知，而在无人机EMAT系统中，从升距到回波信号幅值-相位的完整映射链条涉及电磁场、弹性声场和刚体动力学场的联合求解——其耦合偏微分方程组的解析形式极其复杂且对材料参数（电导率、磁导率、Lam\'{e}常数、表面粗糙度谱）的精确数值高度敏感，难以在工程实践中构建完备且计算可处理的PINN损失项。此外，现有PINN工作几乎全部集中于单一系统的状态估计——如四旋翼的刚体动力学或流体力学——尚未拓展到跨模态、跨时间尺度的异步观测融合问题；基于视觉伺服的方法虽利用相机反馈估计探头-表面相对位姿，但其深度估计精度在近距离（<0.3 m）显著退化，恰与接触建立所需的亚厘米级定位精度构成矛盾。

第二类——物理约束注意力——采取了比PINN更轻量但更针对性的策略：在Transformer的注意力得分矩阵$\mathbf{A} = \mathbf{Q}\mathbf{K}^\top / \sqrt{d}$上直接叠加物理先验偏置$\mathbf{P}$，使修正后的注意力分布$\text{softmax}(\mathbf{A} + \lambda\mathbf{P})$同时受数据驱动力和物理先验力的共同决定\citing{zhao2023pinnsformer}。这一策略的优势在于：物理约束作为加性偏置而非乘性门控，保留了模型在数据充分区域内完全依赖相关性学习的能力，同时在数据稀疏或噪声占优区域内，物理先验可主导注意力分布，使模型输出保持在物理解释可允许的方向上。此外，加性偏置的形式允许根据任务需求灵活设计$\mathbf{P}$的结构——可包含时间邻近矩阵（鼓励相邻帧之间的短程注意力）、空间邻近矩阵（基于六自由度欧氏距离惩罚远距离令牌对的注意力）和法向一致性矩阵（基于姿态角度差调制注意力权重）——而无需修改损失函数或网络架构。

尽管上述两类物理约束策略为解决问题提供了方向，但当前方法存在一个共同的、尚未被充分认识的结构性限制：物理先验被施加于模型的“输出端”（PINN损失函数）或“中间表示层”（物理约束注意力），而非被编码为多模态观测本身的“生成机制约束”。当升距由$h_1$变为$h_2$时，EMAT的期望回波幅值$A(h)$应满足$\|\mathbf{B}(h)\| \propto h^{-n}$的负幂律；当法向速度$v_n$持续负向（即向表面移动）时，接触概率$\hat{p}_t$在物理上不应在相邻帧之间出现大幅下降。这两条约束反映的是传感器数据“应当如何”随着物理状态变化而演化的规范性知识，而非对特定损失函数或网络层的偏好。将这类生成机制约束显式地编码为模型的物理一致性训练目标——即模型输出不仅要拟合标签，还必须与已知物理规律保持定向一致性——是当前方法尚未系统涵盖的方向。

这正是本文逐级物理约束方法论的核心贡献所在。本文不将物理约束视为附加的正则化项或网络架构的修饰，而是将其构建为一条从“多模态观测的物理生成一致性”到“注意力交互的时空合理性”再到“接触概率时序的滞回平滑性”的逐层递进约束链，使物理先验在模型推理的每一个关键环节——数据对齐、特征交互和决策输出——都发挥可辨识、可检验的规范作用。

\section{研究内容与创新}

% 承上：文献综述揭示的四条技术边界 → 四项创新逐一回应

\subsection{主要研究内容}

针对无人机搭载电磁超声探头开展自主探测时面临的接触状态难以可靠判定、多源传感信息难以按物理关系有效融合，以及姿态扰动引起回波质量波动等问题，本文以多旋翼无人机—EMAT探头—被测试件构成的检测系统为研究对象，以物理约束注意力机制为主线，研究从探测状态感知到回波质量反馈调控的方法。具体而言，本文首先研究物理先验对注意力交互的约束及其对接触状态识别的作用，继而拓展至视觉、位姿与电磁超声回波的多模态融合，进一步探讨利用回波特征反馈调节无人机姿态的策略。

\subsubsection{物理约束注意力机制与视觉单模态验证}

标准Transformer的自注意力权重$\text{softmax}(\mathbf{Q}\mathbf{K}^\top/\sqrt{d_k})$衡量的是令牌在嵌入空间中的统计相似性，而非物理世界中时间帧之间的几何连续性、运动约束或接触因果性。在无人机EMAT检测中，被测金属表面的镜面高光、D435深度传感器的近场退化以及PX4悬停漂移三重因素叠加，使得嵌入空间中的视觉相似性与物理接触状态之间产生系统性偏离——Transformer可能将高光饱和像素误判为接触信号。纯数据驱动的注意力分布在训练数据覆盖范围外缺乏物理约束，在临界接触段（接触建立与脱离前后$\pm 2$秒）产生不可解释的误判。

针对上述问题，本文将时间邻近、空间邻近和位姿残差邻近三项物理先验构造为约束矩阵$\mathbf{P} \in \mathbb{R}^{T \times T}$，以对数偏置形式注入自注意力计算：$\text{softmax}(\mathbf{Q}\mathbf{K}^\top/\sqrt{d_k} + \lambda\mathbf{P})$。在视觉单模态的简化实验条件下——仅使用D435深度图的ROI区域提取6维视觉特征（平均深度、深度方差、梯度幅值、归一化深度、像素填充率），排除EMAT回波和位姿残差的融合影响——独立验证物理约束矩阵$\mathbf{P}$对注意力分布和接触分类性能的增益。设计三类消融实验（去除物理约束矩阵$\lambda=0$、去除时序平滑损失、去除物理一致性损失）量化各组件的独立贡献，为第三章后续多模态融合实验提供纯净的理论基线，如图\ref{fig:problem1}所示。

\begin{figure}[htbp]
\centering
\includegraphics[width=0.95\textwidth]{c1/problem1.png}
\caption{问题一及解决方法}
\label{fig:problem1}
\end{figure}

\subsubsection{多模态融合的物理约束注意力与自适应可靠性门控}

EMAT超声回波（约40 Hz）、飞控里程计位姿（约100 Hz）和RGBD视觉深度（约30 Hz）三类传感器的物理生成机制、采样频率和噪声特性存在根本性差异。多模态融合面临两个交织的挑战：其一，各传感器在硬件层面存在不可忽略的固有延迟（视觉曝光约15 ms、EMAT USB批量传输约5 ms、飞控EKF滤波约10 ms），异步观测直接拼接将导致时间错位和虚假跨模态关联；其二，各模态在不同工况下的固有不确定性是非均匀且时变的——视觉在镜面反光或深度空洞时退化，位姿在姿态未收敛时不可靠，超声在升距波动时信噪比骤降——固定权重的等权融合无法在模态退化时自适应降权。

针对上述问题，本文首先建立统一时间基准，以高斯核加权插值将三模态异步观测对齐至约10 ms步长的等间隔时间网格，并通过模态独立编码器（超声一维CNN、位姿MLP、视觉轻量CNN）将异构数据映射至统一隐空间。在此基础上，将物理约束矩阵$\mathbf{P}$的构造参数$(\alpha=0.4, \beta=0.35, \gamma=0.25)$保持不变，从单模态自注意力扩展至多模态交叉注意力——物理约束不仅规范同一模态内相邻令牌的注意力分布，还约束跨模态令牌对之间的交互合理性，抑制EMAT噪声帧与远距离视觉令牌之间的虚假高注意力关联。进一步引入模态可靠性门控向量$\mathbf{g} = [g_Z, g_P, g_V]^\top$，通过可学习的sigmoid门控对各模态实时数据质量进行自适应评估：当视觉出现强反光时$g_V$下降，当姿态未收敛时$g_P$下降，当多模态同时稳定时所有门控权重均接近1。门控融合后的接触概率通过双阈值持续时间滞回状态机（$\tau_{\text{on}}=0.85, \tau_{\text{off}}=0.4$，持续性窗口$K$帧）映射至飞控APPROACH/HOLD状态切换，将连续概率估计转化为具有控制安全保证的离散决策判据，如图\ref{fig:problem2}所示。

\begin{figure}[htbp]
\centering
\includegraphics[width=0.95\textwidth]{c1/problem2.png}
\caption{问题二及解决方法}
\label{fig:problem2}
\end{figure}

\subsubsection{电磁超声回波质量反馈的姿态调控}

第三章的感控网络在估计探测状态的同时输出机体坐标系下的位置与偏航修正量，但宏观控制量本身不能提供改善电磁耦合的姿态微调方向。针对测量过程中微小横滚、俯仰扰动造成的电磁耦合和回波变化，第四章在共享编码器上扩展姿态微调头，将回波质量作为微调分支的触发与评价依据。研究通过能量、信噪比、饱和程度、到达时间及测厚估计的短时稳定性构造可标定的质量评价，由姿态微调头直接回归横滚与俯仰修正量，与第三章宏观控制量叠加并经安全执行层约束后下发。控制器在姿态幅值、变化率、试件法向偏差及位置误差等约束下对微调输出进行门控与投影；回波输入过期或接触状态异常时则冻结微调并退回基准姿态。该反馈链路及实验效果需通过受控姿态扰动、飞控接口验证和闭环测厚试验进一步评估。

\subsection{主要贡献与创新}

针对上述研究现状中揭示的系统性不足：

\begin{enumerate}[
    label=(\arabic*),
    labelindent=\parindent,
    itemindent=0pt,
    labelwidth=2em,
    labelsep=0.5em,
    leftmargin=*
]
\item 力学接触与电磁耦合的二元分离；
\item 多模态融合中统计关联对物理因果性的替代；
\item 动态工况下接触概率的时序闪烁；
\item 感知-控制-检测的串行割裂。
\end{enumerate}

本文以多模态逐级物理约束为核心方法论，以无人机搭载EMAT自主探测为应用场景，开展了从物理建模、平台搭建、信号采集、特征融合到闭环决策的系统性研究。本文在理论方法和系统实现两个层面开展研究，形成了三项主要创新与贡献：

\subsubsection{物理约束注意力机制}

针对Transformer注意力权重衡量统计相似性而非物理因果性的机制缺陷：在自注意力计算中显式引入基于时间邻近和空间残差的物理约束矩阵，以对数偏置形式将无人机运动模型与接触物理先验嵌入跨模态特征交互过程，使注意力分布在数据稀疏或噪声占优区域仍受物理可解释性的规范。

\subsubsection{多模态探测状态感知与控制指令生成框架}

针对力学接触无法推断电磁耦合状态的二元分离问题：首次将EMAT超声回波序列、无人机位姿动态残差及RGBD局部视觉ROI特征进行统一时序建模，将接触状态从瞬时二值标签提升为概率化的连续表达，并由探测状态与控制指令双输出头在共享编码器上端到端生成飞控修正量，使力学状态、电磁耦合状态与飞行决策在同一隐空间中联合表征。

\subsubsection{多源不确定性传播与可靠性门控}

针对动态工况下接触概率时序闪烁导致控制闭环失效的问题：设计了模态可靠性门控权重，对视觉深度失真、位姿未收敛和超声信噪比下降等不同退化模式进行自适应证据分配，并通过双阈值持续时间约束状态机实现从感知概率到控制决策的平滑滞回映射。

\subsubsection{电磁超声回波驱动的姿态反馈调控}

针对接触保持阶段无人机姿态扰动导致电磁耦合与回波质量波动的问题：在多模态感控网络基础上，将连续回波质量特征与位姿状态共同输入共享编码器，由姿态微调头直接回归横滚与俯仰修正量，并在表面法向、位置误差、姿态幅值及变化率等安全约束内经安全执行层投影，形成“回波采集—状态估计—姿态微调—再采集”的闭环。回波能量不作为孤立的最大化目标，而与信噪比、饱和程度及回波/测厚特征的短时稳定性联合评价；具体效果需通过受控姿态扰动实验验证。

\section{本论文的结构安排}

本文共分为五章，各章内容围绕上述研究问题与贡献展开，形成"理论基础—状态感知—反馈调控—实验验证"的递进结构：

第一章：绪论。阐述研究背景与意义，综述非侵入接触式检测无人机、精确接触感控、反馈微调感控和多模态物理信息Transformer四个方向的研究现状，明确本文的切入点与四项核心贡献。

第二章：电磁超声换能与接触检测理论基础。从电磁-弹性耦合机理出发，推导洛伦兹体力与弹性动力学方程的基本关系，建立接触边界与升距效应的回波统计表征模型，为多模态探测状态感知提供物理先验。

第三章：基于多模态融合物理约束注意力机制的无人机电磁超声自适应感控探测方法。先在视觉单模态的简化条件下构造物理约束矩阵$\mathbf{P}$并注入视觉Transformer自注意力，通过消融实验独立验证物理约束与时序平滑损失的增益；再扩展至EMAT超声、飞控位姿和RGBD视觉三模态，研究统一时序对齐、模态独立编码、可靠性门控及探测状态与控制指令双输出头的端到端感控设计，并分析无人机自主贴靠试件与厚度测量的实验结果。

第四章：基于物理约束注意力机制的电磁超声回波反馈无人机姿态调控。在第三章状态估计与控制指令双输出头感控网络基础上，扩展回波质量驱动的横滚/俯仰姿态微调分支，建立带姿态和接触安全约束的微调触发与投影策略，并提出回波质量、姿态稳定性和测厚重复性的闭环评价方案。

第五章：总结与展望。对全文工作进行总结，区分已完成的平台与方法、待实验验证的反馈调控方案，指出当前研究的局限性并展望未来研究方向。
